\chapter{Operators as Values: Constructed Operators and Capability Typing}
\label{ch:operators}

\chapterorient{In \cref{ch:language} we learned that a function is a value, and
that partial application can \emph{freeze} some arguments now and supply the rest
later---``build once, apply many.'' In \cref{ch:strata} we learned that the
build-time configuration a solver bakes in lives in a different stratum, with a
different lifetime, from the run-time vectors it processes. This chapter pays both
ideas off at once. It makes precise the single most important value in the whole
numerical machinery: the \emph{operator}---a function built once with its entire
configuration baked into immutable internal state, then applied again and again as
a pure function. By the end you will be able to say exactly why ``the
preconditioner is just a function you apply,'' why the \emph{same} \code{GMRES}
drives wildly different preconditioners without changing a line, why a \emph{flexible}
preconditioner cannot be a constructed operator and must instead be threaded
through state, and how a phantom-typed \emph{brand} can make ``you passed the
preconditioner where the true operator belongs'' a compile-time error rather than a
silent wrong answer.}

\section{The shape we keep meeting}

We have now met the same shape four or five times, in four or five disguises, and
it is time to name it once and for all.

In \cref{ch:targets} the simulator assembled an operator from the mesh and
materials and then solved a system built from that operator. In \cref{ch:language}
a constructor took a configuration and \emph{returned a function}
(\lstinline[style=l4style]|make_operator :: Config -> (Tensor[N] -> Tensor[N])|),
which we then applied to many different vectors. In \cref{ch:strata} we separated
the \emph{build-time stratum}---configs, tables, a factorization, computed once and
never touched again---from the \emph{run-time stratum}---the iterate, the residual,
the state that evolves every step. Each time, the same picture: an expensive,
once-only \emph{construction} produces a thing; a cheap, repeated \emph{application}
uses it.

That thing is what this book calls a \emph{constructed operator}, and it is the
backbone of every solver, every preconditioner, every assembled matrix in the
chapters ahead. The shape is so pervasive that it earns its own vocabulary, its own
type spelling, and---because getting it slightly wrong silently corrupts a
solve---its own pitfalls. This chapter develops all three.

\begin{keyidea}
A \emph{constructed operator} is a value built \emph{once}, with its entire
configuration baked into immutable internal state, then \emph{applied} many times as
a pure function. Construction is expensive and happens at build time (factorize,
precompute, allocate); application is cheap and pure with respect to the operator's
internals (it reads them, never changes them). The two phases live in two strata
(\cref{ch:strata}) with two lifetimes. Spotting this shape---and resisting the urge
to collapse the two phases---is the central skill of this chapter.
\end{keyidea}

\subsection{Two phases, two strata}

Make the two phases concrete. A constructed operator has exactly two events in its
life, and they can be arbitrarily far apart in the program.

\begin{description}[style=nextline,leftmargin=1.5em]
  \item[Construction.] Inputs are the \emph{configuration}: enum flags, raw tables,
  problem data, a matrix to factorize. The output is an operator value carrying
  internal immutable state. This step is often expensive---a sparse $LU$
  factorization, a precomputed basis-evaluation table, an assembled stiffness
  matrix. It happens \emph{once}.
  \item[Application.] The only input is the \emph{operand}---the vector (or sim
  state) being transformed. The output is a new vector. This step is cheap relative
  to construction, and it is \emph{pure with respect to the operator's internals}:
  it reads the baked-in state, produces a fresh result, and changes nothing. It
  happens \emph{many} times.
\end{description}

\Cref{fig:constructapply} draws the split. Configs and tables flow \emph{in} on the
left; an opaque operator value sits in the middle; a stream of cheap
\code{apply(v)} calls flows out on the right, each reading the same frozen internals
and returning a fresh vector. The left arrow fires once; the right arrows fire
thousands of times.

\begin{figure}[t]
\centering
\begin{tikzpicture}[node distance=10mm]
  % ---- construction inputs ----
  \node[data, align=left, text width=26mm] (cfg)
    {config / enums\\raw tables\\a matrix to\\factorize};
  % ---- the construct step ----
  \node[stage, minimum width=26mm, minimum height=14mm, right=14mm of cfg] (con)
    {\textbf{construct}\\[1pt]\footnotesize factorize,\\precompute};
  \draw[flow] (cfg) -- node[above,font=\scriptsize]{once} (con);
  % ---- the opaque operator value ----
  \node[draw=simBlue, fill=simBgBlue, very thick, rounded corners=3pt,
        minimum width=24mm, minimum height=20mm, right=14mm of con] (op) {};
  \node[font=\ttfamily\scriptsize, simInk] at (op.north) [yshift=-3.2mm]
    {\lstinline[style=l4style]|Op[N -> N]|};
  \node[font=\scriptsize\itshape, simGray, align=center] at (op.center)
    {frozen\\internals\\(factorization,\\tables)};
  \draw[flow] (con) -- (op);
  % ---- many cheap applies out ----
  \foreach \y/\lbl in {1.1/a, 0.0/b, -1.1/c}{
    \node[draw=simTeal, fill=simBgGreen, semithick, rounded corners=2pt,
          minimum width=18mm, minimum height=6mm] (ap\lbl) at ($(op.east)+(2.8,\y)$)
          {\footnotesize\ttfamily apply v};
    \draw[flow] (op.east) -- (ap\lbl.west);
  }
  \node[font=\scriptsize\itshape, simGray] at ($(op.east)+(2.8,-1.9)$) {many times, cheap};
  % build-time / run-time band labels
  \node[font=\scriptsize\bfseries, simRust] at ($(con)+(0,-1.6)$) {build-time stratum};
  \node[font=\scriptsize\bfseries, simGreen] at ($(op.east)+(2.8,1.9)$) {run-time stratum};
\end{tikzpicture}
\caption[The construct-then-apply phase split]{The defining shape of a constructed
operator. \emph{Construction} (left) consumes the configuration---enums, tables, a
matrix to factorize---and runs \emph{once}, producing an opaque operator value whose
internal state (a factorization, precomputed tables) is now frozen. \emph{Application}
(right) consumes only an operand vector and runs \emph{many} times, each call cheap
and pure on the frozen internals. The two phases live in the two strata of
\cref{ch:strata}: construction is build-time, application is run-time, and the
boundary between them is a lifetime boundary the type system can enforce.}
\label{fig:constructapply}
\end{figure}

\subsection{It is currying, made first-class}

You have, in fact, already built constructed operators---you just called them
closures. Recall \code{scaled\_by} from \cref{wex:scaledby}: it took a scalar,
captured it, and returned a function that scales a tensor by that captured value.
Construction was ``supply the scalar''; application was ``supply the tensor''; the
captured scalar was the baked-in internal state. A constructed operator is that
pattern grown up: the captured state is no longer a single scalar but a whole
configuration---a factorization, a stack of basis tables, an enum that selected an
algorithm---and the returned function is no longer a throwaway lambda but a
long-lived value the rest of the program holds and applies.

This is why \cref{ch:language} could already write the signature

\begin{lstlisting}[style=l4style]
make_operator :: Config -> Op[Tensor[N] -> Tensor[N]]
\end{lstlisting}

\noindent and read it as ``give me a configuration, I return an operator value you
apply later.'' \emph{Constructing} the operator is partial application: we supply
the \code{Config} argument now and \emph{freeze} it, getting back a more specialized
value awaiting the run-time vector (\cref{fig:opcurry}). The
\lstinline[style=l4style]|Op[in -> out]| brackets are the same reader-intent marker
as the closure-returning parentheses of \cref{ch:language}---they announce ``the
result is a thing you hold and apply''---with one addition: they promise the held
thing carries \emph{baked-in parameters}, the configuration, as a first-class,
named part of the value rather than as an invisible capture.

\begin{figure}[t]
\centering
\begin{tikzpicture}[node distance=9mm]
  \node[font=\bfseries\small, simInk] at (0,2.5)
    {Construction is partial application: freeze the configuration, await the vector};
  % stage 0: the constructor, needs config + vector
  \node[draw=simTeal, fill=simBgGray, thick, rounded corners=2pt,
        minimum width=24mm, minimum height=11mm] (f0) at (-4.6,1.0) {};
  \node[font=\ttfamily\scriptsize] at (-4.6,1.32) {make\_operator};
  \node[font=\scriptsize, simTeal] at (-4.6,0.78) {needs cfg, $v$};
  \node[font=\scriptsize\itshape,simGray,below=0.5mm of f0,align=center]
    {\code{Config ->}\\\code{Op[N -> N]}};
  % supply config
  \draw[flow] (-3.3,1.0) -- node[above,font=\scriptsize]{supply \code{cfg}} (-1.6,1.0);
  % stage 1: the constructed operator, config frozen
  \node[draw=simBlue, fill=simBgBlue, thick, rounded corners=2pt,
        minimum width=24mm, minimum height=11mm] (f1) at (-0.2,1.0) {};
  \node[font=\ttfamily\scriptsize] at (-0.2,1.32) {apply\_A = make\_operator cfg};
  \node[font=\scriptsize, simBlue] at (-0.2,0.78) {needs $v$};
  \node[font=\scriptsize\itshape,simGray,below=0.5mm of f1] {\code{Op[N -> N]}};
  \node[draw=simRust,fill=simBgRust,rounded corners=1pt,font=\scriptsize,
        inner sep=1.5pt] at (-0.2,2.05) {\code{cfg} frozen in};
  % supply v
  \draw[flow] (1.0,1.0) -- node[above,font=\scriptsize]{supply $v$} (2.6,1.0);
  % stage 2: a vector
  \node[draw=simRust, fill=simBgRust, thick, rounded corners=2pt,
        minimum width=16mm, minimum height=11mm] (r) at (3.8,1.0) {$\mat{A}v$};
  \node[font=\scriptsize\itshape,simGray,below=0.5mm of r] {\code{Tensor[N]}};
\end{tikzpicture}
\caption[Constructing an operator is partial application]{A constructed operator is
currying (\cref{ch:language}) raised to first class. \code{make\_operator} needs a
configuration \emph{and}, later, a vector. Supplying the configuration does not
produce a vector---it produces a \emph{new value}, the constructed operator
\code{apply\_A}, with the configuration \emph{frozen in}. Supplying a vector to that
operator finally produces a result. The two arrows are the two phases of
\cref{fig:constructapply}; the frozen \code{cfg} is the operator's immutable internal
state. The \lstinline[style=l4style]|Op[...]| spelling differs from a plain closure
only in advertising that the frozen part is a first-class, named configuration.}
\label{fig:opcurry}
\end{figure}

\subsection{Application is pure on the internals}

The word ``pure'' is doing real work in ``applied as a pure function,'' and it is
worth pinning down, because it is exactly the property that makes constructed
operators safe to share.

Application reads the operator's internal state and never writes it. So two
applications of the same operator cannot interfere: \code{apply\_A x1} and
\code{apply\_A x2} each read the frozen factorization, each produce a fresh result,
and neither disturbs the other or the operator. This is the immutability discipline
of \cref{ch:language} applied to the operator's internals: there is nothing to clone
and nothing to protect, because the internals cannot be overwritten by any
application. An operator, once constructed, is as inert and shareable as the number
$7$.

\begin{pitfall}
``Pure on the internals'' is not the same as ``stateless.'' A constructed operator
\emph{does} carry state---a whole factorization, perhaps megabytes of it. The
discipline is that this state is \emph{set at construction and never changed
thereafter}. If you find yourself wanting an operator to update something between
applications---an iteration count, a running residual, a cached previous input---stop:
that thing is \emph{not} operator-internal state. It is sim state (\cref{ch:strata}),
and it must be threaded explicitly through the run-time stratum, not stuffed into the
operator. An operator that mutates between applications has stopped being a pure
function, and every guarantee in this chapter evaporates with it. We return to
exactly this failure---and why it is not a flaw but a real distinction---in
\cref{sec:decisionrule}.
\end{pitfall}

\section{Worked example: a constructed Jacobi preconditioner}

The cleanest possible constructed operator is the \emph{Jacobi} (diagonal)
preconditioner, and tracing it end to end makes every abstract claim above concrete.
We will meet it again, in its proper numerical setting, in \cref{ch:precond}; here it
is purely a vehicle for the construct-then-apply shape.

\begin{workedexample}{Building and applying a Jacobi preconditioner}{jacobi}
A preconditioner is an \emph{approximate inverse}: given a matrix $\mat{A}$, we want
a cheap operator $\mat{M}^{-1}$ with $\mat{M}^{-1}\approx \mat{A}^{-1}$, so that
applying it to a residual $\vr$ yields something close to the error. The Jacobi
choice is the cheapest possible: take $\mat{M}=\operatorname{diag}(\mat{A})$, the
diagonal of $\mat{A}$, and ignore everything off the diagonal. Then $\mat{M}^{-1}$ is
just ``divide each component by the corresponding diagonal entry.''

\textbf{Construction (once).} Extract the diagonal of $\mat{A}$ and invert it
componentwise, producing a single vector $\mathbf{d}$ of reciprocals. This is the entire
baked-in internal state:
\begin{lstlisting}[style=l4style]
make_jacobi :: Tensor[N, N] -> Op[Tensor[N] -> Tensor[N]]
make_jacobi A =
  let diag = extract_diagonal A in   -- read the N diagonal entries: O(N)
  let d    = reciprocal diag in      -- d[i] = 1 / A[i,i]: O(N)
  \r -> d .* r                       -- the closure: captures d, awaits r
\end{lstlisting}
The returned closure captures \code{d}---the operator's frozen internal state---and
awaits a residual. Construction touched $\mat{A}$ exactly twice (extract, invert),
each an $O(N)$ pass; it produced one length-$N$ vector. Notice what construction did
\emph{not} do: it never looked at any residual, because no residual exists yet.

\textbf{Application (many times).} Inside a solve, every iteration produces a fresh
residual $\vr_k$ and applies the preconditioner to it:
\begin{lstlisting}[style=l4style]
let pc = make_jacobi A in       -- build ONCE, before the loop
... inside the iteration ...
let z1 = pc r1 in               -- z1 = d .* r1 ; O(N), pure on d
let z2 = pc r2 in               -- z2 = d .* r2 ; O(N), reads the SAME d
let z3 = pc r3 in               -- ...
\end{lstlisting}
Each application is one componentwise multiply: $\vz_k = \mathbf{d} \odot \vr_k$, an $O(N)$
operation. Take $N=4$, $\operatorname{diag}(\mat{A})=[4,2,8,1]$, so $\mathbf{d} =
[\tfrac14,\tfrac12,\tfrac18,1]$. Then for $\vr_1=[8,6,16,3]$,
\[
  \vz_1 \;=\; \mathbf{d} \odot \vr_1 \;=\; \bigl[\tfrac14\!\cdot\!8,\ \tfrac12\!\cdot\!6,\
  \tfrac18\!\cdot\!16,\ 1\!\cdot\!3\bigr] \;=\; [2,\,3,\,2,\,3],
\]
and for a \emph{different} residual $\vr_2=[4,4,4,4]$ on the \emph{same} operator,
\[
  \vz_2 \;=\; \mathbf{d} \odot \vr_2 \;=\; \bigl[1,\,2,\,\tfrac12,\,4\bigr].
\]

\textbf{Three readings, each a habit from \cref{ch:language}.} (1) \emph{The phases
are separate and asymmetric in cost.} Construction did two $O(N)$ passes over the
matrix, once. Each application does one $O(N)$ multiply, but there are thousands of
them; the matrix is never re-examined. The reciprocal $\mathbf{d}$ was computed once and
amortized over the whole solve. (2) \emph{Application is pure on the captured
\code{d}.} Computing $\vz_1$ did not change $\mathbf{d}$; $\vz_2$ read exactly the same
$\mathbf{d}$; the two could be computed in either order, or in parallel, without affecting
each other (\cref{fig:exprtree}'s ``no hidden order,'' now applied to operator
application). (3) \emph{Nothing was overwritten.} \code{r1}, \code{r2}, and \code{d}
all name exactly the values they did before each call; each \code{z} is a brand-new
vector.

The Jacobi preconditioner is the constructed-operator pattern in its purest form: the
expensive-but-once part (form $\mathbf{d}$) is frozen into the operator, and the cheap part
(multiply by $\mathbf{d}$) is repeated. Every richer preconditioner in
\cref{ch:precond,ch:multigrid}---an incomplete factorization, a multigrid V-cycle---is
the \emph{same shape} with more elaborate baked-in state.
\end{workedexample}

\section{The decision rule: construction-bound versus per-step variants}
\label{sec:decisionrule}

We now reach the load-bearing distinction of the chapter---the one rule that, applied
correctly, tells you whether a piece of configuration belongs inside a constructed
operator or out in the threaded state. Getting it right is the difference between a
clean solve and a subtle, convergent, \emph{wrong} one.

\begin{keyidea}
A variant that is \emph{bound once at construction} and held fixed for the whole
solve $\to$ a constructed operator absorbs it cleanly. A variant that \emph{changes
per step}---a different value on each iteration---\emph{cannot} be a constructed
operator, because there is no single value to freeze at construction. The per-step
variant must instead enter the threaded \code{SimState} (\cref{ch:strata,ch:monad}).
This single rule is exactly why \emph{flexible} preconditioners exist, and exactly
why \code{FGMRES} is a different algorithm from \code{GMRES}.
\end{keyidea}

The reasoning is almost a tautology once stated. Construction happens \emph{once},
before the loop. To freeze a value at construction, that value must \emph{exist}, and
be \emph{final}, at that moment. A preconditioner $\mat{M}$ chosen at solve start and
held fixed satisfies this: it exists, it is final, freeze it. But a preconditioner
$\mat{M}_k$ that is allowed to be \emph{different on each step $k$}---a so-called
\emph{flexible} preconditioner---does not: at construction time the later $\mat{M}_k$
do not yet exist, and there is no single $\mat{M}$ to bake in. The construction phase
has nothing to capture.

When this happens, the variant has not vanished---it has merely moved. A value that
changes every step \emph{is} run-time state, by definition. So it belongs in the
run-time stratum: it must be threaded through \code{SimState}, evolving step by step
like the iterate and the residual, rather than frozen in the operator. The
construction-vs-run-time stratification of \cref{ch:strata} is precisely the test:
\emph{ask which stratum the variant's lifetime puts it in.}

\begin{figure}[t]
\centering
\begin{tikzpicture}[node distance=8mm]
  % decision diamond
  \node[draw=simGold, fill=simBgGold, thick, diamond, aspect=2.2, inner sep=1pt,
        align=center, font=\footnotesize] (q) at (0,2.6)
    {Does the variant\\change \emph{per step}?};
  % NO branch (left): constructed operator
  \node[stage, minimum width=30mm, align=center, below left=12mm and 2mm of q] (con)
    {\footnotesize bind \emph{once} at construction\\[1pt]\textbf{constructed operator}};
  \node[draw=simBlue, fill=simBgBlue, semithick, rounded corners=2pt,
        minimum width=34mm, align=center, font=\scriptsize, below=5mm of con] (conex)
    {build $\mat{M}^{-1}$ once;\\\code{apply(pc, r)} each step;\\\code{op} never changes};
  \draw[flow] (con) -- (conex);
  \draw[flow] (q.south west) to[out=210,in=90]
    node[left,font=\scriptsize\bfseries,simGreen]{NO} (con.north);
  % YES branch (right): threaded state
  \node[product, minimum width=30mm, align=center, below right=12mm and 2mm of q] (thr)
    {\footnotesize thread it through state\\[1pt]\textbf{SimState schema grows}};
  \node[draw=simRust, fill=simBgRust, semithick, rounded corners=2pt,
        minimum width=34mm, align=center, font=\scriptsize, below=5mm of thr] (threx)
    {add per-step $\mat{M}_k$ data\\to the carry (\code{Z} basis);\\evolves every iteration};
  \draw[flow] (thr) -- (threx);
  \draw[flow] (q.south east) to[out=-30,in=90]
    node[right,font=\scriptsize\bfseries,simRust]{YES} (thr.north);
  % captions
  \node[font=\scriptsize\itshape, simGray, below=4mm of conex, text width=36mm, align=center]
    {fixed preconditioner;\\GMRES, CG};
  \node[font=\scriptsize\itshape, simGray, below=4mm of threx, text width=36mm, align=center]
    {flexible preconditioner;\\FGMRES};
\end{tikzpicture}
\caption[The decision rule]{The one rule that places a variant. Ask: does it change
per step? If \emph{no} (a preconditioner chosen at solve start and held fixed), bind
it once at construction---a constructed operator absorbs it, and the iteration calls
\code{apply(pc, r)} uniformly. If \emph{yes} (a flexible preconditioner $\mat{M}_k$
differing each step), no construction can freeze it; the variant must enter the
threaded \code{SimState} and evolve every iteration. The right branch is why
\code{FGMRES} carries an extra basis the left branch does not---the subject of
\cref{wex:flexible} and, in full, \cref{ch:gmres}.}
\label{fig:decisionrule}
\end{figure}

\subsection{Worked example: fixed versus flexible, in two signatures}

The decision rule is easiest to feel as a contrast between two signatures. We put the
fixed preconditioner and the flexible one side by side and watch the variant move out
of the operator and into the state.

\begin{workedexample}{Why a flexible preconditioner cannot be a constructed
operator}{flexible}
A Krylov solver (\cref{ch:gmres}) builds, step by step, a basis $\{\vv_0, \vv_1,
\dots\}$ of search directions, applying the preconditioner to each. We compare two
ways the preconditioner can enter.

\textbf{Case 1: a fixed preconditioner (a constructed operator works).} The
preconditioner $\mat{M}$ is chosen once, at solve start. We construct it before the
loop and apply it---the \emph{same} operator---at every step:
\begin{lstlisting}[style=l4style]
-- the preconditioner is a constructed operator, built once:
let pc = make_preconditioner config A in   -- pc :: Op[Tensor[N] -> Tensor[N]]
-- each Arnoldi step applies the SAME pc to a fresh basis vector:
let z_j = pc v_j in                        -- z_j = M^{-1} v_j, same M every j
let w   = matvec A z_j in ...
\end{lstlisting}
Here \code{pc} appears once, at construction, and is then applied uniformly. The state
the loop threads is the basis $V = [\vv_0, \dots, \vv_j]$ and some bookkeeping; the
preconditioner is \emph{not} part of that state, because it never changes. The
constructed operator absorbs the entire ``which preconditioner'' question. This is
the left branch of \cref{fig:decisionrule}.

\textbf{Case 2: a flexible preconditioner (it must move into the state).} Now the
preconditioner is allowed to differ each step---$\mat{M}_j$ at step $j$ (perhaps it is
itself an inner iterative solve whose accuracy varies, perhaps it adapts to the
residual). There is no single $\mat{M}$ to construct. The application becomes
\begin{lstlisting}[style=l4style]
-- there is NO single pc to construct; M_j is produced per step:
let z_j = apply_M_step j v_j in   -- z_j = M_j^{-1} v_j, a DIFFERENT M_j each j
\end{lstlisting}
and---this is the crucial consequence---the per-step preconditioned vectors $\vz_j$
can no longer be recomputed on demand from a frozen operator, because the operator
that made them is gone by the next step. They must be \emph{remembered}: the solver
threads a \emph{second} basis $Z = [\vz_0, \dots, \vz_j]$ alongside $V$. Contrast the
two state schemas:
\begin{lstlisting}[style=l4style]
-- GMRES  (fixed pc): the preconditioner is a constructed operator, NOT in state
type GmresState  = { x: Vec, V: Vec[], H: Mat, ... }
-- FGMRES (flexible pc): the per-step preconditioned basis Z IS threaded state
type FgmresState = { x: Vec, V: Vec[], Z: Vec[], H: Mat, ... }
--                                     ^^^^^^^ the extra basis the flexible case forces
\end{lstlisting}
The solution update reflects the move exactly: GMRES recovers the answer from the
plain basis, $\vx = \vx_0 + V_m \vy_m$, while FGMRES must use the remembered
preconditioned basis, $\vx = \vx_0 + Z_m \vy_m$. The variant that the constructed
operator absorbed in Case~1 has, in Case~2, become a field of the threaded state.

\textbf{The reading.} Nothing went wrong; a real distinction surfaced. A fixed
preconditioner is build-time configuration---freeze it, the operator is a value. A
flexible preconditioner is run-time state---thread it, the carry grows by one basis.
The decision rule (\cref{fig:decisionrule}) reads the variant's \emph{lifetime} and
routes it to the matching stratum. ``\code{FGMRES} exists because the preconditioner
became per-step state'' is the whole story, and \cref{ch:gmres} is where we tell it in
full.
\end{workedexample}

\begin{pitfall}
The tempting mistake is to try to keep a flexible preconditioner inside a constructed
operator by letting the operator \emph{mutate}---``I'll just update \code{pc}'s
internal $\mat{M}$ each step.'' This re-introduces exactly the mutable box
\cref{ch:language} spent a whole section banishing. A constructed operator that
mutates between applications is no longer pure on its internals, no longer safely
shareable, and no longer has a well-defined ``value'' at all---``what does \code{pc}
hold \emph{now}?'' becomes a live question again. The discipline is firm: per-step
variation goes in the threaded state, never in a secretly-mutating operator. The cost
is one honest extra field in the state schema; the alternative is a silent corruption
of the value model.
\end{pitfall}

\section{Solver as operator: an inverse is a function you apply}

We have been calling a preconditioner ``an approximate inverse'' and applying it like
an operator. It is time to make the type-level move that justifies this explicit,
because it is one of the most powerful abstractions in the whole numerical stack.

\begin{keyidea}
A \emph{solver} is an operator. A solver computes (or approximates) the action
$\vv \mapsto \mat{M}^{-1}\vv$ of an inverse, and that action has \emph{exactly the
same shape} as the action $\vv \mapsto \mat{M}\vv$ of the operator it inverts: vector
in, vector out. So a solver can wear the same \code{apply} interface as a plain
operator, and be substituted anywhere an operator is expected. ``Apply the
preconditioner to the residual'' is just \code{apply(pc, r)}---a function call,
indistinguishable in form from \code{apply(A, x)}.
\end{keyidea}

This sounds almost too simple to be worth a name, but its consequence is enormous. A
Krylov method (\cref{ch:gmres}) is written against the \code{apply} interface and
nothing else. It asks the operator $\mat{A}$ for matrix--vector products
\code{apply(A, v)}, and it asks the preconditioner for $\code{apply(pc, r)} \approx
\mat{M}^{-1}\vr$---and that is \emph{all} it ever asks of either. The iteration never
sees the \emph{algorithm} inside the preconditioner. It does not know, and cannot
ask, whether \code{apply(pc, r)} runs a diagonal multiply, an incomplete
factorization, a multigrid V-cycle, or a direct sparse solve. It sees a black box
that turns a residual into an approximate error, and it iterates.

\begin{figure}[t]
\centering
\begin{tikzpicture}[node distance=10mm]
  % the GMRES driver
  \node[stage, minimum width=26mm, minimum height=22mm] (gm) at (0,0)
    {\textbf{GMRES}\\[2pt]\footnotesize asks only:\\\code{apply(A, v)}\\\code{apply(pc, r)}};
  % the operator black box (top)
  \node[draw=simBlue, fill=simBgBlue, thick, rounded corners=2pt,
        minimum width=22mm, minimum height=9mm, right=24mm of gm, yshift=8mm] (A)
        {\footnotesize operator $\mat{A}$};
  % the preconditioner black box (bottom) -- SAME interface
  \node[draw=simViolet, fill=simBgViolet, thick, rounded corners=2pt, densely dashed,
        minimum width=22mm, minimum height=9mm, right=24mm of gm, yshift=-8mm] (pc)
        {\footnotesize solver \code{pc}};
  % the SAME apply arrow to both
  \draw[flow] (gm.east|-A) -- node[above,font=\scriptsize]{\code{apply}} (A.west);
  \draw[flow] (gm.east|-pc) -- node[below,font=\scriptsize]{\code{apply}} (pc.west);
  % what's INSIDE pc, hidden
  \node[align=left, font=\scriptsize\itshape, simGray, text width=36mm, right=4mm of pc]
    {inside (hidden): Jacobi?\\ILU? multigrid V-cycle?\\direct solve?\\\textbf{GMRES cannot tell.}};
  \node[align=left, font=\scriptsize\itshape, simGray, text width=30mm, right=4mm of A]
    {a true operator action\\$v \mapsto \mat{A}v$};
  % brace: both wear the same interface
  \draw[decorate, decoration={brace, amplitude=4pt}, simGray]
    ($(A.north west)+(-0.2,0.2)$) -- ($(pc.south west)+(-0.2,-0.2)$)
    node[midway, left=5pt, font=\scriptsize\bfseries, simInk, align=center]
    {same\\\code{apply}\\interface};
\end{tikzpicture}
\caption[A solver wears the operator's interface]{The solver-as-operator move. The
\code{GMRES} driver is written against one interface---\code{apply}---and asks for two
kinds of action through it: matrix--vector products from the true operator $\mat{A}$,
and approximate-inverse actions from the preconditioner \code{pc}. Because a solver
\emph{is} an operator (same vector-in, vector-out shape), \code{pc} wears the very
same \code{apply} interface as $\mat{A}$, and the iteration cannot see---and never asks
about---the algorithm hidden inside it. This is exactly why the same \code{GMRES}
drives a Jacobi preconditioner, an incomplete factorization, and a multigrid V-cycle
without changing a line: the differences are all behind the black box.}
\label{fig:solveroperator}
\end{figure}

\Cref{fig:solveroperator} draws the picture. The payoff is the one the whole book
keeps returning to: \emph{write the coordination once, feed it the verbs}
(\cref{ch:language}). Here the coordination is the Krylov iteration; the verbs are
\code{apply(A, $\cdot$)} and \code{apply(pc, $\cdot$)}; and because both verbs have the
same type, the \emph{same} iteration drives every combination of operator and
preconditioner the simulator offers. The driven solver, the electrostatic solver, and
the eigenmode solver's inner solves all reuse one \code{GMRES}---they differ only in
which operator and which preconditioner they hand it, and those are just two
\code{apply}-shaped values.

\subsection{Why this is a rotation, not a renaming}

It is fair to ask whether ``a solver is an operator'' says anything at all, or merely
relabels one thing as another. It says something, and the test is whether any
\emph{state} is hidden by the move.

At the source level (\cref{ch:strata} called this $L_0$), a solver is genuinely a
different kind of thing from a bare operator: it carries extra machinery---a routine to
set which operator it inverts, an initial-guess slot, convergence tolerances, internal
iteration state. The solver-as-operator rotation \emph{hides all of that} from the
iteration. The Krylov loop sees only the \code{apply} action; the solver's
operator-setting routine, its tolerances, its guess slot are compressed away,
invisible behind the interface. Because real state is hidden---because the loop
\emph{cannot} reach the preconditioner's internal apparatus even though it is
there---the move is a genuine \emph{rotation} (\cref{ch:rotations}), not a typedef.

The contrast sharpens it. If ``solver is an operator'' meant only that we wrote
\code{Solver} as an alias for \code{Operator}, with no extra structure underneath,
then the statement would hide nothing and buy nothing. It is precisely because the
solver \emph{does} carry distinct apparatus---which the iteration is then prevented
from touching---that wearing the operator interface is an act of abstraction. The
chapter on rotations (\cref{ch:rotations}) makes this ``does it hide state?'' test the
formal criterion for an honest rotation.

\section{The constructed-operator factory: where the variant is consumed}

So far each constructed operator has had one configuration. But a real solver must
choose \emph{among} configurations: which Krylov method (\code{CG} / \code{GMRES} /
\code{FGMRES}), which preconditioner (Jacobi / AMS / algebraic multigrid / a direct
factorization), which side, which orthogonalization. These choices arrive as
\emph{enums}---discrete flags read from a config file. Where do the enums get
consumed?

The answer is the \emph{constructed-operator factory}: a single function that takes
the enum (and any context it needs), dispatches on it internally, and returns a
\emph{uniformly typed} operator. The factory is the one place the variant is read.
Before it, the code branches on ``which method?''; after it, the code holds one
operator value and applies it---variant-free.

\begin{figure}[t]
\centering
\begin{tikzpicture}[node distance=9mm]
  % the enum coming in
  \node[data, align=center, text width=24mm] (enum)
    {variant enum\\\code{CG | GMRES |}\\\code{FGMRES | ...}};
  % the factory
  \node[stage, minimum width=34mm, minimum height=16mm, right=12mm of enum] (fac)
    {\textbf{factory}\\[1pt]\footnotesize \code{ConfigureKrylovSolver}\\\footnotesize dispatch \emph{here, once}};
  \draw[flow] (enum) -- (fac);
  % a uniformly-typed operator out
  \node[draw=simBlue, fill=simBgBlue, very thick, rounded corners=2pt,
        minimum width=24mm, minimum height=12mm, right=14mm of fac] (op)
        {\footnotesize\ttfamily Op[N -> N]};
  \draw[flow] (fac) -- node[above,font=\scriptsize]{one type} (op);
  % variant-free code after
  \node[draw=simGreen, fill=simBgGreen, thick, rounded corners=2pt,
        minimum width=30mm, minimum height=12mm, right=12mm of op, align=center,
        font=\footnotesize] (after)
        {downstream code:\\\code{apply(op, v)}\\\textbf{no enum in sight}};
  \draw[flow] (op) -- (after);
  % the "wall" after which variants are gone
  \draw[simRust, very thick, densely dashed] ($(fac.east)+(0.7,1.7)$) -- ($(fac.east)+(0.7,-2.0)$);
  \node[font=\scriptsize\bfseries, simRust, align=center] at ($(fac.east)+(0.7,-2.4)$)
    {variant\\wall};
  \node[font=\scriptsize\itshape, simGray] at ($(enum)+(0,-2.0)$) {variant-bearing};
  \node[font=\scriptsize\itshape, simGray] at ($(after)+(0,-2.0)$) {variant-free};
\end{tikzpicture}
\caption[The factory: a variant enum in, a uniform operator out]{The
constructed-operator factory is the single site at which a variant enum is consumed.
To its left, code carries the enum and must branch on it; the factory dispatches on
the enum \emph{once} and returns an operator of one uniform type. To its right---past
the ``variant wall''---downstream code holds that single operator and applies it with
\code{apply(op, v)}, never re-inspecting the enum. Palace's \code{ConfigureKrylovSolver}
(Krylov method) and \code{ConfigurePreconditionerSolver} (preconditioner type) are exactly
this: enum in, uniform solver-operator out, variant-free code after.}
\label{fig:factory}
\end{figure}

\Cref{fig:factory} draws the factory as a ``variant wall.'' Palace has two of them,
structurally identical: \code{ConfigureKrylovSolver} consumes the
Krylov-method$\,\times\,$side$\,\times\,$orthogonalization$\,\times\,$restart axes and
returns an iterative-solver operator; \code{ConfigurePreconditionerSolver} consumes the
preconditioner-type$\,\times\,$multigrid$\,\times\,$scalar-field axes and returns a
preconditioner operator. Both read a clutch of enums, dispatch once, and hand back a
single uniformly typed value the rest of the solver applies without ever asking which
method it got. A solver's signature schematically:

\begin{lstlisting}[style=l4style]
-- the factory: ONE place the krylov-method enum is read
configure_krylov :: KrylovMethod -> Config -> Op[Tensor[N] -> Tensor[N]]
-- downstream, the method is gone; only the uniform operator remains:
let ksp = configure_krylov method cfg in   -- dispatch happens HERE, once
let x   = apply ksp b in                    -- ... and never again
\end{lstlisting}

This is the procedural heart of \emph{variant absorption}, which the next section
makes into a discipline. For now hold the single observation: the factory localizes
the entire ``which variant?'' decision to one call. Everywhere else in the solver, an
operator is just an operator.

\section{Capability typing: distinguishing roles without runtime cost}

A factory hands back ``a uniformly typed operator.'' But uniformity can go too far.
Consider Palace's preconditioned Krylov solver: it holds \emph{two} operators of the
\emph{same} underlying type \code{Op[N -> N]}---the \emph{true operator} $\mat{A}$ it is
solving against, and the \emph{preconditioner-assembly operator} from which the
preconditioner is built. Both are \code{Op[N -> N]}. Nothing in their type stops a
careless caller from passing them in the wrong order. And if they \emph{are} swapped,
the solve does not crash---it runs, and converges, to the \emph{wrong answer}. A wrong
solve that converges is the worst kind of bug: silent.

\emph{Capability typing} is the discipline that turns this silent runtime bug into a
loud compile-time type error, at zero runtime cost.

\begin{keyidea}
A \emph{capability type} is a \emph{phantom-typed brand}: a marker attached to a value
that records the value's intended \emph{role}, enforced by the type checker but with
no runtime representation---no extra field, no allocation, no dispatch. Two values can
share an underlying type yet carry distinct brands, so that using one where the other
is required is a \emph{type error} rather than a silently-wrong computation. The brand
costs nothing at runtime and is erased before the program runs; its entire job is to
make a class of role-confusion mistakes \emph{unspeakable}.
\end{keyidea}

\subsection{The brand, concretely}

Concretely, we take the underlying operator type and attach a phantom tag that
distinguishes the role:

\begin{lstlisting}[style=l4style]
-- underlying type, shared:
type Op[N -> N] = ...
-- two BRANDED views, distinct at the type level, identical at runtime:
type TrueOp       = Op[N -> N] & { __cap: "true" }         -- the operator we solve against
type PcAssemblyOp = Op[N -> N] & { __cap: "pc_assembly" }  -- the operator the pc is built from
-- smart constructors: runtime identities, type-level commitments
as_true_op        :: Op[N -> N] -> TrueOp
as_pc_assembly_op :: Op[N -> N] -> PcAssemblyOp
\end{lstlisting}

\noindent The tag \code{\_\_cap} is \emph{phantom}: there is no such field at runtime,
no byte of storage, no branch. The smart constructors \code{as\_true\_op} and
\code{as\_pc\_assembly\_op} are the identity function at runtime---they return their
argument untouched. Their \emph{only} job is to commit a value to a role at the moment
of construction, so that a downstream signature can demand a specific brand:

\begin{lstlisting}[style=l4style]
set_operators :: TrueOp -> PcAssemblyOp -> Solver -> Solver
\end{lstlisting}

\noindent Now passing a \code{PcAssemblyOp} where a \code{TrueOp} is required is
rejected by the type checker---a compile-time error, caught before any number is
computed---even though the two values are byte-for-byte identical at runtime. The
permitted ``escape hatch'' is intended and useful: when you genuinely want to
precondition the true operator against itself, you brand the \emph{one} operator
\emph{both} ways (\code{as\_true\_op A} and \code{as\_pc\_assembly\_op A}); the
discipline forbids passing one brand where the other is wanted, not applying both
brands to one value.

\begin{figure}[t]
\centering
\begin{tikzpicture}[node distance=8mm]
  % LEFT: variant absorption -- hides axes of variation
  \begin{scope}
    \node[font=\bfseries\small, simInk] at (0,2.7) {Variant absorption};
    \node[font=\scriptsize\itshape, simGray] at (0,2.3) {hides axes of \emph{variation}};
    % the uniform operator, with variants collapsed inside
    \node[draw=simBlue, fill=simBgBlue, thick, rounded corners=2pt,
          minimum width=30mm, minimum height=22mm] (u) at (0,0.6) {};
    \node[font=\ttfamily\scriptsize] at (u.north) [yshift=-3mm] {Op[N -> N]};
    \node[font=\scriptsize, simGray, align=center] at (u.center)
      {\code{CG}? \code{GMRES}?\\\code{Jacobi}? \code{AMG}?\\\emph{collapsed inside,}\\\emph{one type out}};
    \node[font=\scriptsize\itshape, simGray, below=2mm of u, text width=40mm, align=center]
      {many configurations $\to$ \emph{one} uniform type};
  \end{scope}
  % RIGHT: capability typing -- distinguishes axes of role
  \begin{scope}[shift={(7.4,0)}]
    \node[font=\bfseries\small, simInk] at (0,2.7) {Capability typing};
    \node[font=\scriptsize\itshape, simGray] at (0,2.3) {distinguishes axes of \emph{role}};
    % two branded operators, same underlying type, kept apart
    \node[draw=simTeal, fill=simBgGreen, thick, rounded corners=2pt,
          minimum width=26mm, minimum height=8mm] (tt) at (0,1.25) {};
    \node[font=\ttfamily\scriptsize] at (tt.center) {TrueOp};
    \node[draw=simRust, fill=simBgRust, thick, rounded corners=2pt,
          minimum width=26mm, minimum height=8mm] (pp) at (0,-0.1) {};
    \node[font=\ttfamily\scriptsize] at (pp.center) {PcAssemblyOp};
    % they share an underlying type (a thin bracket)
    \draw[decorate, decoration={brace, amplitude=3pt, mirror}, simGray]
      ($(tt.east)+(0.15,0)$) -- ($(pp.east)+(0.15,0)$)
      node[midway, right=4pt, font=\scriptsize, simGray, align=left]
      {same\\underlying\\\code{Op[N -> N]}};
    % the "no swap" barrier between them
    \node[font=\scriptsize\bfseries, simRust] at (0,0.58) {swap = type error};
    \node[font=\scriptsize\itshape, simGray, below=4mm of pp, text width=40mm, align=center]
      {one type $\to$ \emph{two} brands kept apart};
  \end{scope}
\end{tikzpicture}
\caption[Capability brands versus variant axes]{The two disciplines are
complementary, pulling in opposite directions on the same operator type.
\emph{Left:} variant absorption \emph{hides} axes of \emph{variation}---\code{CG} vs
\code{GMRES}, Jacobi vs multigrid---collapsing many configurations behind one uniform
\code{Op[N -> N]}. \emph{Right:} capability typing \emph{distinguishes} axes of
\emph{role}---\code{TrueOp} vs \code{PcAssemblyOp}---splitting one underlying type into
two branded views the checker refuses to interchange. Absorption makes ``which
configuration?'' unspeakable downstream; branding makes ``which role?'' unswappable. A
single operator can be both absorbed (its variant hidden) and branded (its role
fixed).}
\label{fig:capvariant}
\end{figure}

\subsection{Background: where brands come from}

The phantom-brand idea is borrowed, lightly, from typed functional programming, and
you may meet it under several names. In Haskell, a \code{newtype} wraps an existing
type in a distinct, zero-cost name the compiler keeps separate---the same
representation, a different type. In Rust, a \emph{zero-sized marker type} (a field of
a type with no data, such as \code{PhantomData}) tags a value with a role that exists
only for the borrow checker and vanishes from the compiled binary. The word
\emph{capability} comes from the security literature, where a capability is an
unforgeable token of authority; here the ``authority'' is the role---holding a
\code{TrueOp} is the authority to act as the operator the iteration solves against,
holding a \code{PcAssemblyOp} is the authority to act as the source of the
preconditioner. We use brands in their type-checking-only form: they constrain what
the program can express, and then they are erased. \Citet{peytonjones2003} is the
standard reference for the underlying type machinery.

\subsection{Capability typing versus variant absorption}

Capability typing and variant absorption are easy to confuse because both are about
``many operators of the same type,'' but they pull in opposite directions, and seeing
the contrast cleanly is worth the effort (\cref{fig:capvariant}).

\begin{itemize}[nosep]
  \item \textbf{Variant absorption hides axes of \emph{variation}.} \code{CG} versus
  \code{GMRES}, Jacobi versus multigrid, left- versus right-side: these are
  \emph{differences we want to disappear} behind a uniform interface, so downstream
  code need not branch on them. Absorption \emph{collapses} many things into one type.
  \item \textbf{Capability typing distinguishes axes of \emph{role}.} True operator
  versus preconditioner-assembly operator: these are \emph{differences we want to
  preserve}, because confusing them is a bug. Branding \emph{splits} one type into
  several that the checker refuses to interchange.
\end{itemize}

A single operator can be \emph{both} at once. The preconditioner-assembly operator is
\emph{absorbed} along its variation axes---it is uniform across real versus complex
scalar field, uniform across which preconditioner algorithm sits inside it---and
simultaneously \emph{branded} \code{PcAssemblyOp}, distinct from \code{TrueOp} along
the role axis. Absorption answers ``which configuration is this?'' with ``don't ask,
it's uniform''; branding answers ``which role does this play?'' with ``that, and the
checker will hold you to it.'' They compose without conflict because they speak about
different axes of the same value.

\section{Variant absorption as a discipline}

We have now used ``absorb the variant'' informally several times. It is time to make
it a discipline with a checkable standard, because the failure mode---absorbing a
variant \emph{partly}, while pretending to absorb it fully---is one of the most
insidious ways a specification goes quietly wrong.

The principle: when a piece of the machinery has \emph{orthogonal axes of
variation}---preconditioner side, orthogonalization kind, fixed-versus-flexible,
real-versus-complex scalar field---the form should absorb those axes
\emph{parametrically}, as bindings of a parameter, rather than \emph{appending} them
as special-case paragraphs bolted onto a base description. ``Here is \code{GMRES};
oh, and for \code{FGMRES} replace this with that; oh, and for left-preconditioning add
this step'' is the anti-pattern. The parametric form names a parameter once and lets
its binding range over the variants.

\subsection{The three levels of absorption}

Absorption is not all-or-nothing; it holds at three nested levels, and a genuine
absorption holds at \emph{all three} (\cref{fig:absorptionlevels}).

\begin{description}[style=nextline,leftmargin=1.5em]
  \item[(a) Invariant level: the math unifies.] A single equation or recurrence covers
  every variant, perhaps with the variant appearing as a parameter. Example: write the
  solution update as $\vx = \vx_0 + W_m \vy_m$ with $W_m = V_m$ for \code{GMRES} and
  $W_m = Z_m$ for \code{FGMRES}. One equation; the parameter $W$ selects the variant.
  This is necessary but not sufficient.
  \item[(b) Procedural level: the variant is mentioned once.] The procedure names the
  variant at most once---at the binding or dispatch site (the factory call)---and never
  re-inspects it. ``Let \code{op = configure(side); then each step apply(op, v)}'' is
  procedurally absorbed: \code{side} appears once, at construction. A procedure that
  re-tests the variant at the operator-application step \emph{and} the update step is
  \emph{not} procedurally absorbed, even if its math unifies.
  \item[(c) Primitive-sequence level: the operation chain has the same shape.] The
  underlying sequence of primitive calls---which operations, in which order---is
  identical across variants, with the variant binding only the \emph{operands}, not the
  chain. ``$[\text{matvec},\,\text{dot},\,\text{axpy}]$ for every \code{side}, with the
  matvec's operator argument changing'' is primitive-sequence absorbed. A variant that
  forces one branch to run an extra primitive the others skip is \emph{not}.
\end{description}

\begin{figure}[t]
\centering
\begin{tikzpicture}[node distance=4mm]
  % three nested bands, (a) outermost
  \node[band, fill=simBgGold, draw=simGold, thick, minimum width=82mm, minimum height=46mm]
    (a) at (0,0) {};
  \node[font=\footnotesize\bfseries, simGold!75!black, anchor=north] at (a.north) [yshift=-1.5mm]
    {(a) invariant level --- one equation covers all variants};
  \node[band, fill=simBgBlue, draw=simBlue, thick, minimum width=66mm, minimum height=32mm]
    (b) at (0,-2.5mm) {};
  \node[font=\footnotesize\bfseries, simBlue, anchor=north] at (b.north) [yshift=-1.5mm]
    {(b) procedural --- variant named once};
  \node[band, fill=simBgGreen, draw=simGreen, thick, minimum width=50mm, minimum height=18mm]
    (c) at (0,-5mm) {};
  \node[font=\footnotesize\bfseries, simGreen, anchor=north] at (c.north) [yshift=-1.5mm]
    {(c) primitive-sequence};
  \node[font=\scriptsize, simInk, align=center] at (c.center) [yshift=-1mm]
    {same op chain,\\operands differ};
  % full absorption note
  \node[font=\scriptsize\itshape, simGray, anchor=north, text width=74mm, align=center]
    at (a.south) [yshift=-3mm]
    {\textbf{Full absorption} = all three nested levels hold. The dangerous case is
    reaching (a) but silently \emph{not} (b) or (c)---the form \emph{looks} uniform but
    the operation chain secretly branches.};
\end{tikzpicture}
\caption[The three nested levels of variant absorption]{Absorption holds at three
nested levels. \emph{(a)} The mathematical statement unifies into one
parameter-bearing equation. \emph{(b)} The procedure mentions the variant at most
once, at its binding site, and never re-inspects it. \emph{(c)} The underlying
primitive-call sequence has the same shape across variants, the variant binding only
the operands. Full absorption requires all three. The insidious failure is partial,
\emph{undisclosed} absorption---reaching (a) while quietly violating (b) or (c)---which
the next pitfall names.}
\label{fig:absorptionlevels}
\end{figure}

\subsection{The anti-pattern: silent partial absorption}

The three levels exist so that we can name precisely the way absorption goes wrong.

\begin{pitfall}
\textbf{Silent partial absorption} is the failure this entire discipline exists to
catch. A form achieves invariant-level absorption (a)---one clean equation, looks
uniform---but \emph{silently} fails procedural (b) or primitive-sequence (c)
absorption: under the surface, the procedure re-inspects the variant at several sites,
or the operation chain branches by variant, and the form's tidy appearance hides it.
The reader sees ``$\vx = \vx_0 + W_m \vy_m$, one equation'' and believes the variant
is handled; the actual computation forks three ways. The fix is honesty: if (b) or (c)
genuinely cannot hold---because the variant is structural, like a flexible
preconditioner forcing an extra threaded basis (\cref{wex:flexible})---then
\emph{disclose} the residual divergence explicitly, listing the sites where the
procedure or the primitive chain forks. A disclosed partial absorption is fine; a
\emph{silent} one is a specification that lies about its own uniformity.
\end{pitfall}

The routes to genuine full absorption are the constructions of this chapter. A
\emph{constructed operator} that internalizes the variant at construction time
achieves (b) and (c) at once: the factory reads the enum once, the operator is
applied uniformly, the primitive chain is one shape (\cref{fig:factory}). When even a
constructed operator cannot help---because the variant is per-step, and the decision
rule (\cref{fig:decisionrule}) routes it into the threaded state---then the absorption
is genuinely partial, and the discipline demands you say so, in the open, rather than
let a tidy equation imply a uniformity the operation chain does not have.

\begin{remark}
Notice how the pieces of this chapter fit together as one mechanism. The
\emph{factory} (\cref{fig:factory}) is \emph{where} a variant is absorbed---the single
site that reads the enum. The \emph{constructed operator} (\cref{fig:constructapply})
is \emph{what} absorbs it---the value that internalizes the choice and is then applied
uniformly. The \emph{decision rule} (\cref{fig:decisionrule}) tells you \emph{whether}
absorption can succeed at all---no, if the variant is per-step. And \emph{capability
typing} (\cref{fig:capvariant}) handles the orthogonal concern that survives
absorption---the \emph{roles} the uniform operators play, which must stay distinct even
after their variations are hidden. Absorption hides what varies; branding preserves
what must not be confused.
\end{remark}

\section{Putting it together: the preconditioned solver}

We can now read the whole construction surface of a real preconditioned Krylov solver
as one coherent picture, with every piece of this chapter in its place. Palace's
\code{BaseKspSolver} does exactly the following, and nothing in it should now be
mysterious.

At \emph{construction} (build-time stratum, run once): two factory calls
(\cref{fig:factory}). \code{ConfigureKrylovSolver} consumes the
Krylov-method$\,\times\,$side$\,\times\,$orthogonalization$\,\times\,$restart enums and
returns a constructed iterative-solver operator \code{ksp}; \code{ConfigurePreconditionerSolver}
consumes the preconditioner-type$\,\times\,$multigrid$\,\times\,$scalar-field enums and
returns a constructed preconditioner operator \code{pc}. Both are \code{apply}-shaped
(solver-as-operator, \cref{fig:solveroperator}); both have absorbed their variation
axes parametrically inside the factory; after the two calls, no enum survives into the
solve. The two operators are then \emph{branded}---the true operator \code{as\_true\_op},
the preconditioner-source \code{as\_pc\_assembly\_op} (\cref{fig:capvariant})---so the
binding step \code{set\_operators :: TrueOp -> PcAssemblyOp -> ...} cannot receive them
swapped.

At \emph{application} (run-time stratum, run many times): the solve folds the Krylov
iteration to convergence, and at each step it asks the two constructed operators for
their \code{apply} actions and nothing else---matrix--vector products from \code{ksp}'s
operator, approximate-inverse actions from \code{pc}. The iteration never re-examines
any enum, never sees either operator's internal algorithm, and never confuses their
roles, because the construction phase settled all three concerns once and froze the
result. The cleanly separated, frozen construction is exactly what makes the cheap,
uniform, repeated application possible.

This is the shape every solver in \cref{ch:gmres,ch:precond,ch:multigrid} wears. The Krylov methods of
\cref{ch:gmres}, the preconditioners of \cref{ch:precond}, and the multigrid V-cycle
of \cref{ch:multigrid} are all elaborations of this one surface: build constructed
operators by reading enums in a factory; brand their roles; thread only the genuinely
per-step state; apply the operators uniformly through a single \code{apply}
interface. Operators as values is not one technique among many---it is the frame the
entire numerical machinery hangs on.

\summarysection
A \emph{constructed operator} is a value built once, with its full configuration baked
into immutable internal state, then applied many times as a pure function. The two
phases live in two strata (\cref{ch:strata}): construction is expensive, build-time,
and run once; application is cheap, run-time, run many times, and pure on the frozen
internals. This is the first-class form of the currying and ``build once, apply many''
of \cref{ch:language}; the \lstinline[style=l4style]|Op[in -> out]| spelling marks a
held-and-applied value carrying named baked-in parameters. The \emph{decision rule} is
load-bearing: a variant bound \emph{once at construction} (a fixed preconditioner) is
absorbed by a constructed operator; a variant that changes \emph{per step} (a flexible
preconditioner $\mat{M}_k$) cannot be---there is no single value to freeze---so it must
enter the threaded \code{SimState} instead. This is exactly why \code{FGMRES} threads
an extra basis \code{Z} that \code{GMRES} does not (\cref{ch:gmres}). A \emph{solver is
an operator}: an approximate inverse wears the same \code{apply} interface as the thing
it inverts, so a Krylov method needs only the black-box action $\vv \mapsto
\mat{M}^{-1}\vv$ and never sees the preconditioner's algorithm---which is why one
\code{GMRES} drives wildly different preconditioners. The
\emph{constructed-operator factory} is the single site a variant enum is consumed and a
uniformly typed operator returned (\code{ConfigureKrylovSolver},
\code{ConfigurePreconditionerSolver}); after it, the code is variant-free.
\emph{Capability typing} attaches a phantom-typed brand recording a value's \emph{role}
(\code{TrueOp} vs \code{PcAssemblyOp}) at zero runtime cost, turning a role-swap into a
type error rather than a silent wrong-but-convergent solve; it is complementary to
\emph{variant absorption}, which hides axes of \emph{variation} where branding
distinguishes axes of \emph{role}. Variant absorption is a checkable discipline with
three nested levels---(a) the math unifies, (b) the variant is named once, (c) the
primitive chain has one shape---with full absorption requiring all three and
\emph{silent partial absorption} the anti-pattern it exists to catch.

\furtherreading
The constructed-operator and solver-as-operator patterns are standard in operator-algebra
preconditioning frameworks; \citet{saad2003} (especially the chapter on preconditioned
Krylov methods) motivates the central abstraction---that a Krylov method needs only the
black-box action $\vv \mapsto \mat{M}^{-1}\vv$, so any approximation respecting that
interface composes uniformly---and is the reference for the flexible-preconditioner
distinction that drives \code{FGMRES}. The phantom-type and zero-cost-brand machinery
behind capability typing is part of the type-system toolkit of typed functional
programming; \citet{peytonjones2003} is the standard reference for the underlying
mechanisms (\code{newtype} wrappers and phantom type parameters), which Rust mirrors
with zero-sized marker types. The two disciplines---hiding variation and preserving
role---meet again in \cref{ch:rotations}, where ``does the move hide real state?''
becomes the formal test that separates an honest rotation from a mere renaming.

\exercisessection
\begin{enumerate}[nosep]
  \item \textbf{Two phases, two costs.} For the constructed Jacobi preconditioner of
  \cref{wex:jacobi}, state the cost of \emph{construction} and the cost of a
  \emph{single application}, each in terms of $N$. If a solve runs $k$ iterations,
  write the total work as a function of $N$ and $k$, and identify which term is the
  amortized construction. For what range of $k$ is constructing the operator clearly
  worthwhile, and when would calling the inner multiply directly be just as good?
  \item \textbf{Apply the rule.} For each variant below, decide whether it belongs in a
  constructed operator or in the threaded \code{SimState}, and justify your answer with
  the decision rule (\cref{fig:decisionrule}): (a) the time-step size $\Delta t$ in a
  \emph{fixed}-step time integrator; (b) the preconditioner in a flexible Krylov
  method, where an inner solve of varying accuracy is used each step; (c) the matrix
  entries of an assembled stiffness operator, fixed for the whole solve; (d) an
  adaptive damping parameter recomputed from the residual at each iteration.
  \item \textbf{The mutating-operator trap.} A colleague proposes to support a flexible
  preconditioner by giving the constructed \code{pc} an \emph{update} method that
  overwrites its internal $\mat{M}$ each step, ``so we don't have to change the state
  schema.'' Explain, using the value model of \cref{ch:language}, exactly which
  guarantee this breaks, and what question becomes ill-posed once \code{pc} can mutate.
  What is the honest cost of doing it correctly instead?
  \item \textbf{Same interface, different insides.} The same \code{GMRES} drives a
  Jacobi preconditioner, an incomplete factorization, and a multigrid V-cycle without
  changing a line. Using \cref{fig:solveroperator}, name the single interface that
  makes this possible, state precisely what \code{GMRES} asks of the preconditioner,
  and explain what it is \emph{prevented} from asking. Why is being prevented from
  seeing the preconditioner's algorithm a feature, not a limitation?
  \item \textbf{Rotation or renaming?} Recall the test from \cref{ch:rotations} (and
  the discussion in this chapter): a move is a genuine rotation only if it \emph{hides
  real state}. The solver-as-operator move makes a \code{Solver} usable wherever an
  \code{Operator} is expected. List two pieces of apparatus a real solver carries that
  a bare operator does not, and explain why hiding them from the iteration makes this a
  rotation rather than a typedef. What would have to be true for it to be \emph{only} a
  renaming?
  \item \textbf{Where the enum dies.} A solver reads a config flag
  \code{method $\in$ \{CG, GMRES, FGMRES\}}. In a well-absorbed design, at how many sites
  in the per-step iteration does \code{method} appear, and where is the one place it is
  consumed? Sketch (in the notation) the factory call and one downstream apply, and mark
  the ``variant wall'' of \cref{fig:factory}. Then describe what a \emph{silent partial
  absorption} of this enum would look like and why a critic should reject it.
  \item \textbf{Brand the binding.} Two operators \code{A} and \code{P}, both of type
  \code{Op[N -> N]}, are to be passed to
  \lstinline[style=l4style]|set_operators :: TrueOp -> PcAssemblyOp -> Solver -> Solver|.
  (i) Write the two smart-constructor calls that brand them correctly. (ii) Explain why
  swapping the arguments is now a compile-time error even though \code{A} and \code{P}
  are byte-for-byte the same kind of value at runtime. (iii) Show how to legitimately
  precondition \code{A} against itself, and explain why the brand discipline permits
  this while still forbidding the accidental swap. (iv) What is the runtime cost of all
  this branding?
  \item \textbf{Absorption versus branding.} In one or two sentences each, state what
  \emph{variant absorption} does and what \emph{capability typing} does, and identify
  which axis of an operator each one speaks about. Then give a single operator that is
  simultaneously absorbed along one axis and branded along another, and name both axes
  explicitly (\cref{fig:capvariant}).
  \item \textbf{The three levels, by hand.} Consider preconditioner side
  $\in \{\text{LEFT},\text{RIGHT},\text{NONE}\}$ in a Krylov step. Suppose a draft
  writes the step as one equation (level (a)) but, under the surface, the LEFT branch
  computes \lstinline[style=l4style]|w = apply(pc, apply(A, v))|, the RIGHT branch
  \lstinline[style=l4style]|w = apply(A, apply(pc, v))|, and NONE just
  \lstinline[style=l4style]|w = apply(A, v)|. Which of levels (b) and (c) does this draft
  fail, and why? Describe the constructed operator that would repair both at once, and
  state what its construction freezes.
\end{enumerate}
